\section{Keystone-refinement tests}
\label{sec:keystone_refinement}

The machinery of
Sections~\ref{sec:relaxation_protocol}--\ref{sec:convergence}
addresses one failure mode of a long-timeline restriction: the
restriction is \emph{obsolete}---the pathway no longer activates when
exercised, so the restriction can be lifted entirely.  In practice a
second failure mode is more common: the restriction is
\emph{misconfigured}---the pathway still activates and lifting the
restriction wholesale would remain dangerous, but technological
progress has made narrower keystones available that would block the
risk pathway without blocking adjacent benign pathways.  This section
extends the controlled relaxation construction to handle the
misconfiguration case.

\subsection{Keystone structures and their partial order}

A restriction $\Xirestrict$ is realised by restricting specific
keystone capabilities in the tuple that would otherwise enable the
risk pathway~\citep[Proposition~1]{lasser2026horizon}.  Different
choices of which capabilities to restrict may block the same
pathway.  We make this explicit.

\begin{definition}[Keystone structure]
\label{def:keystone_structure}
Given a risk claim
$R = (a_j, (d_1, \ldots, d_k), \Pathway, \pR)$, a
\emph{keystone structure} for $R$ is a subset
$\keystonestruct \subseteq \{d_1, \ldots, d_k\}$ whose restriction
blocks the exercise pathway $\Pathway$: if every $d \in
\keystonestruct$ is restricted, $\Pathway$ is not exercisable.  Write
$\keystonestruct \vdash \Pathway$ when this holds.
\end{definition}

\begin{definition}[Keystone partial order]
\label{def:keystone_order}
Given two keystone structures $\keystonestruct_1, \keystonestruct_2$
that both block $\Pathway$, write
$\keystonestruct_1 \preceq \keystonestruct_2$ iff
$\keystonestruct_1 \subseteq \keystonestruct_2$.  The relation
$\preceq$ is a partial order on blocking keystone structures; a
$\preceq$-minimal element is a \emph{minimal blocking keystone
structure}.  Restriction cost is monotone along $\preceq$: narrower
structures impose less restriction on adjacent capabilities.
\end{definition}

\begin{remark}[Technology drift]
\label{rem:technology_drift}
The set of candidate keystone structures evolves with technology.
A keystone that was $\preceq$-minimal at time $t_0$ may become
non-minimal at $t_1 > t_0$ if new capability differentiation
(for instance, new manufacturing tolerances, new detection
methods, or new isolation protocols) produces a finer keystone
$\keystonestruct' \prec \keystonestruct$ that still blocks the
pathway.  The framework's channel for surfacing such finer
structures is the structural discovery
value~\citep[Proposition~2]{lasser2026horizon}: if structural
discovery reveals a new capability that enables a narrower
blocking structure, the discovery is value-positive regardless of
whether the refinement is subsequently accepted.
\end{remark}

\subsection{Refinement tests as controlled relaxation}

A keystone-refinement test probes whether a candidate narrower
structure $\keystonestruct_{\mathrm{narrow}} \prec
\keystonestruct_{\mathrm{current}}$ still blocks $\Pathway$.

\begin{definition}[Keystone-refinement test]
\label{def:refinement_test}
A \emph{keystone-refinement test} for risk claim $R$ with current
restriction realised by $\keystonestruct_{\mathrm{current}}$ is a
controlled relaxation test in the slots of
Definition~\ref{def:cr_test} with the third slot taken to be the
restriction \emph{being lifted}:
$T = (R, \Xirestrict_{\mathrm{lifted}}, \scope, \testdur,
\damagebound)$, where
$\Xirestrict_{\mathrm{lifted}}$ is the restriction induced by the
keystone difference
$\keystonestruct_{\mathrm{current}} \setminus
\keystonestruct_{\mathrm{narrow}}$ for some candidate
$\keystonestruct_{\mathrm{narrow}} \prec
\keystonestruct_{\mathrm{current}}$.  The narrower structure
$\keystonestruct_{\mathrm{narrow}}$ remains in force throughout
the test as ambient restriction during $\scope$; the test asks
whether $\keystonestruct_{\mathrm{narrow}} \vdash \Pathway$ holds
empirically by exercising only the released keystones.  The
scope-construction and damage-bound conditions of
Definition~\ref{def:cr_test} apply unchanged with
$\Xirestrict := \Xirestrict_{\mathrm{lifted}}$.
\end{definition}

The restriction is not fully lifted; only the keystones outside the
narrower structure are released.  The risk pathway must therefore
route through $\keystonestruct_{\mathrm{narrow}}$ to activate, and
$\keystonestruct_{\mathrm{narrow}}$ remains blocking.  Under the
scope and damage apparatus already proved for CR, the refinement
test inherits the damage bound and convergence theorems.

\begin{corollary}[Keystone-refinement under CR damage bound]
\label{cor:refinement_damage}
A keystone-refinement test satisfies the damage bound of
Theorem~\ref{thm:damage_bound}: worst-case
$\volL$-contraction during the test is $\leq \damagebound$ and
post-reinstatement is $\leq 2 \damagebound$.  The test is a
specialisation of Definition~\ref{def:cr_test} with
$\Xirestrict := \Xirestrict_{\mathrm{lifted}}$
(Definition~\ref{def:refinement_test}); $\maxcon(\scope)$ is
computed over pathways that the released keystones permit,
holding the unreleased $\keystonestruct_{\mathrm{narrow}}$
restricted as ambient.  The scope-bound construction
(Definition~\ref{def:test_scope}) and the cooperative-closure
argument (Theorem~\ref{thm:damage_bound}(c)) apply without
modification.
\end{corollary}

\begin{corollary}[Convergence of refinement tests]
\label{cor:refinement_convergence}
Under the operating assumptions of
Theorem~\ref{thm:convergence} (full-exercise regime), repeated
keystone-refinement tests of
$\keystonestruct_{\mathrm{narrow}}$ converge almost surely to the
correct binary hypothesis: $H_0$ that
$\keystonestruct_{\mathrm{narrow}} \vdash \Pathway$ (the narrower
structure blocks the pathway) versus $H_1$ that it does not.
Convergence rate is the KL-divergence rate of
Theorem~\ref{thm:convergence}(d) applied to the refinement-test
observation model.
\end{corollary}

\begin{remark}[Two failure modes of a restriction]
\label{rem:two_failure_modes}
A restriction can be wrong in two ways:
\begin{itemize}
\item[\textbf{Mode A (Obsolete).}]  The pathway no longer activates
  at all; the risk has disappeared from the current context.
  Detected by controlled relaxation of the full restriction
  (Theorem~\ref{thm:convergence}): $\pR$ converges to the
  $H_0$-consistent value with $L_0 < 1$ per test.
\item[\textbf{Mode B (Misconfigured).}]  The pathway still
  activates if the restriction is lifted wholesale, but
  technological differentiation has made a narrower blocking
  structure $\keystonestruct_{\mathrm{narrow}} \prec
  \keystonestruct_{\mathrm{current}}$ available.  Detected by
  keystone-refinement testing
  (Corollary~\ref{cor:refinement_convergence}): the current
  framework finds that $\keystonestruct_{\mathrm{narrow}}$
  suffices and the difference keystones can be released.
\end{itemize}
Both modes are addressed by the same CR apparatus; the distinction
is in which restriction the test probes (the full
$\Xirestrict_{\mathrm{current}}$ for Mode A, the reduced
$\Xirestrict_{\mathrm{narrow}}$ for Mode B).  A Mode A failure does not preclude Mode B success against a
narrower candidate: a restriction that cannot be fully lifted
may still admit refinement.
\end{remark}

\begin{remark}[Meta-policy for refinement]
\label{rem:refinement_policy}
The test-initiation criteria of
Section~\ref{sec:meta_policy} extend to refinement.  The
information value of a refinement test is the expected reduction
in restriction cost from adopting
$\keystonestruct_{\mathrm{narrow}}$ if the test confirms it, minus
the test's expected damage.  When a structural argument is
available that $\keystonestruct_{\mathrm{narrow}} \vdash \Pathway$
(for instance, a mechanical proof that the released keystones
cannot participate in $\Pathway$), the structural argument is
strictly preferable~\citep[Proposition~2]{lasser2026horizon}:
refinement testing should be reserved for cases where the
structural argument is unavailable or contested.
\end{remark}

\subsection{Worked example: nuclear dual-use keystone refinement}
\label{sec:nuclear_example}

The nuclear scenario from~\citep[Definition~2 and
discussion]{lasser2026horizon} illustrates Mode B in practice.  The
capability tuple is $(d_1, d_2, d_3) = (\text{nuclear physics},
\text{enrichment}, \text{weapons machining})$.  The exercise
pathway $\Pathway = $ ``build and detonate a fission weapon''
leads to large $\volL$-contraction proportional to the number of
agents destroyed.

\paragraph{Historical keystone choice.}
Mid-20th-century technology did not operationally distinguish
reactor-grade enrichment (3--5\% fissile content, using standard
centrifuge cascades) from weapons-grade enrichment (90\%+, using
the same cascade technology reconfigured for higher stage count).
A centrifuge cascade capable of the former could in principle be
reconfigured for the latter; separation at the level of individual
keystones was not feasible.  The broad keystone structure
\[
\keystonestruct_{\mathrm{current}}
   = \{d_1, d_2, d_3\}
\]
(restrict all three capabilities broadly) was adopted and blocks
$\Pathway$ at maximum restriction cost.  Under the technology of
the time, $\keystonestruct_{\mathrm{current}}$ was
$\preceq$-minimal.

\paragraph{Modern keystone availability.}
Subsequent technological differentiation has exposed narrower
keystones.  Modern isotope-separation technology distinguishes
cascade configurations by separative work units (SWU) and stage
count in ways that make weapons-grade enrichment identifiable
independently of reactor-grade enrichment.  Weapons machining
(specific tolerance classes for implosion lens fabrication) is
likewise separable from general machining.  Modern candidate
structures include
\[
\keystonestruct_{\mathrm{narrow}}
   = \{d_2^{\mathrm{weapons\text{-}grade}},
       d_3^{\mathrm{weapons\text{-}specific}}\}
\]
which restricts weapons-grade enrichment cascades and
weapons-specific machining tolerances while permitting
reactor-grade enrichment and general machining.  Under modern
technology, $\keystonestruct_{\mathrm{narrow}} \prec
\keystonestruct_{\mathrm{current}}$ if $\keystonestruct_{\mathrm{narrow}}
\vdash \Pathway$.

\paragraph{Refinement test.}
A keystone-refinement test probes whether
$\keystonestruct_{\mathrm{narrow}}$ suffices.  The scope $\scope$
consists of a monitored civilian reactor development programme,
exercising the reactor-pathway (reactor-grade enrichment, general
machining) under $\keystonestruct_{\mathrm{narrow}}$; weapons-grade
enrichment and weapons-specific machining remain restricted.  The
detection channel monitors for indicators of weapons-pathway
activation: enrichment to weapons-grade fissile content, diversion
of reactor feedstock toward weapons-specific cascades, acquisition
of weapons-specific machining tolerances.  Test duration is set
by the time constant of the detection architecture on these
indicators; damage tolerance $\damagebound$ bounds the realised
contraction even under a diversion attempt bounded by the scope.

\paragraph{Outcomes.}
\begin{itemize}
\item If $\keystonestruct_{\mathrm{narrow}} \vdash \Pathway$
  (weapons pathway is not exercisable under
  $\keystonestruct_{\mathrm{narrow}}$, modern technology
  hypothesis): the test observes no weapons-pathway activation
  over the test window, $\pR$ converges toward the consistent
  hypothesis, and a refined restriction regime adopting
  $\keystonestruct_{\mathrm{narrow}}$ becomes operationally
  justified.  Civilian reactor development proceeds without
  collateral restriction.
\item If $\keystonestruct_{\mathrm{narrow}} \not\vdash \Pathway$
  (the weapons pathway remains exercisable through some route
  the narrower structure fails to block---for example, covert
  diversion of reactor-grade enrichment toward weapons use via
  routes modern detection would catch but not always): the test
  observes weapons-pathway activation within $\scope$, the
  safety cutoff fires, damage is bounded by $2 \damagebound$
  (Theorem~\ref{thm:damage_bound}(b)), and
  $\keystonestruct_{\mathrm{current}}$ remains in force.  The
  refinement test itself is the evidence that broader
  restriction remains needed under current technology.
\end{itemize}

\paragraph{Contrast with Mode A testing.}
A Mode A test on the nuclear pathway would ask ``can
$\keystonestruct_{\mathrm{current}}$ be lifted entirely?''  The
pathway still operates; the Mode A test would detect weapons-pathway
activation and correctly reject lifting the restriction.  But the
Mode A finding does not address the operationally more interesting
question: whether the \emph{scope} of restriction can be refined.
Mode B testing is the mechanism for that refinement.  The two tests
are complementary; the choice is not between them but between
whether to perform one, the other, or (in practice) both under a
programmatic review schedule.

\paragraph{Applicability beyond nuclear.}
The keystone-refinement pattern applies wherever a long-timeline
restriction was adopted under technology that did not distinguish
benign from risky keystones.  Environmental-contamination controls
(persistent chemicals whose modern reformulations operate via
different mechanisms), export controls on dual-use technology
(where modern differentiation allows narrower controls), biomedical
restrictions on gain-of-function research (where specific
biocontainment tiers separate research protocols), and other
regulatory domains exhibit the same structure---the policy
literature on dual-use research review~\citep{tucker2012innovation}
collects extensive case studies of exactly this dynamic across
biological and chemical capabilities.  In each, the framework's
contribution is to formalise the refinement test as a
bounded-damage empirical probe rather than a
moratorium-vs.-wholesale-liberalisation binary, and to provide
formal damage-bound and convergence guarantees that the
case-by-case review processes operate without.
